%% LyX 2.4.4 created this file.  For more info, see https://www.lyx.org/.
%% Do not edit unless you really know what you are doing.
\documentclass[english]{article}
\usepackage[T1]{fontenc}
\usepackage[latin9]{inputenc}
\usepackage{enumitem}
\usepackage{amsmath}
\usepackage{amsthm}
\usepackage{geometry}
\geometry{verbose,tmargin=1in,bmargin=1in,lmargin=1in,rmargin=1in}
\PassOptionsToPackage{normalem}{ulem}
\usepackage{ulem}

\makeatletter
%%%%%%%%%%%%%%%%%%%%%%%%%%%%%% Textclass specific LaTeX commands.
\numberwithin{equation}{section}
\numberwithin{figure}{section}
\newlength{\lyxlabelwidth}      % auxiliary length 

%%%%%%%%%%%%%%%%%%%%%%%%%%%%%% User specified LaTeX commands.
\usepackage{fancyhdr}
\usepackage{lastpage}
\pagestyle{fancy}
\usepackage{enumitem}
\usepackage{ifthen}
%sets math fonts to be more readable
\everymath=\expandafter{\the\everymath\displaystyle}
%\DeclareMathSizes{10}{10}{10}{10}
\usepackage{multicol}
% Added by lyx2lyx
\setlength{\parskip}{\smallskipamount}
\setlength{\parindent}{0pt}

\makeatother

\usepackage{babel}
\begin{document}
\renewcommand{\author}{Dr.\,James Ashe}

\newcommand{\classNum}{331}

\newcommand{\classSec}{A}

\newcommand{\nameType}{Name:}

\newcommand{\examType}{Exam 4}

\renewcommand{\date}{January 7, 2013} %%\today

%%First page's header%%%%%%%%%%%%%%%%%%%%%%
\fancypagestyle{firststyle} %%header settings for first pagr
{
	\fancyhead[L]{MTH \classNum{}(\classSec{}) \author{}\\
	\textbf{\examType{}} \date{}}
	\fancyhead[C]{\textbf{\nameType}}
	\setlength{\headsep}{14pt}
}
%%%%%%%%%%%%%%%%%%%%%%%%%%%%%%%%%%%%%%%%%%

%%Next page's header%%%%%%%%%%%%%%%%%%%%%%%
\setlist{nolistsep}
\lhead{MTH \classNum{}(\classSec{})} 
\chead{\textbf{\nameType}} 
\cfoot{\thepage{}~of~\pageref{LastPage}} 
\setlength{\headsep}{5pt} 
%%%%%%%%%%%%%%%%%%%%%%%%%%%%%%%%%%%%%%%%%%%

%%Various settings; see the preamble for other settings%%%%%
%%\setenumerate[0]{leftmargin=12pt,itemindent=0pt,labelwidth=0pt}
\renewcommand{\labelenumi}{\textbf{\arabic{enumi}.}}
\renewcommand{\labelenumii}{\textbf{\alph{enumii}.}}
\thispagestyle{firststyle}
%%%%%%%%%%%%%%%%%%%%%%%%%%%%%%%%%%%%%%%%%%

Unless otherwise stated, each problem is worth 4 points. Show all
of your work.

\emph{Evaluate the following limits.}
\begin{enumerate}[resume]
\item ${\displaystyle \lim_{(x,y)\rightarrow(1,\pi)}(3x+\sin xy})$\vfill
\item ${\displaystyle \lim_{(x,y)\rightarrow(1,-1)}\frac{y^{3}+y^{2}x}{y+x}}$\vfill
\end{enumerate}
\emph{Use the Two-Path Test to prove that the following limit do not
exist.}
\begin{enumerate}[resume]
\item ${\displaystyle \lim_{(x,y)\rightarrow(0,0)}\frac{x^{2}-4y^{3}}{x^{2}+2y^{3}}}$\vfill
\end{enumerate}
\emph{Find the indicated partial derivative of the following functions.}
\begin{enumerate}[resume]
\item Find $g_{x}$ and $g_{y}$ for $g(x,y)=(x^{2}+2)e^{y}$.\vfill\newpage
\item {[}8 pts{]} Find the four second partial derivatives of $h(x,y)=x^{2}\cos2y$.\vspace{3cm}\vfill
\end{enumerate}
\emph{Use the Chain Rule to find the following derivative. When feasible
express your answer in terms of the independent variable. }
\begin{enumerate}[resume]
\item {[}6 pts{]} $dz/dt$ where $z=x^{3}y-x$, $x=t^{2}$, and $y=2t^{-1}$.\vspace{1cm}\vfill
\end{enumerate}
\emph{Compute the gradient of the following function and evaluate
it at the given point P.}
\begin{enumerate}[resume]
\item $f(x,y)=x^{2}+y^{2}$; $P(-1,2)$\vfill\newpage
\end{enumerate}
\emph{Consider the following functions and points P. }

\textbf{a.}\emph{ Find the }\emph{\uline{unit}}\emph{ vectors that
give the direction of steepest ascent and steepest descent at P.}

\textbf{b.}\emph{ Find a vector that points in a direction of no
change in the function at P.}
\begin{enumerate}[resume]
\item {[}6 pts{]} $f(x,y)=x^{2}+y^{2}$; $P(-1,2)$\vfill
\end{enumerate}
\emph{Use the Second Derivative Test to determine (if possible) whether
each critical point of the given function is a local maximum, local
minimum, or saddle point.}
\begin{enumerate}[resume]
\item {[}8 pts{]} $f(x,y)=4+x^{3}+y^{3}-3xy$ with critical points $(0,0)$
and $(1,1).$\vfill
\end{enumerate}
\emph{Find all of the critical points (there are two) of the following
function.}
\begin{enumerate}[resume]
\item $f(x,y)=\frac{x^{3}}{3}-\frac{y^{3}}{3}+3xy$\vfill
\end{enumerate}

\end{document}
